\section{Principales aportaciones}

Para explicar las aportaciones, las dividiremos por hipótesis.

\subsection{Hipótesis \HOne}

Podemos ver el número de parámetros para la representación del estado global o ver cada clase de manera individual. Si nos fijamos en la ilustración \ref{fig:experimentsByNInputs}, podemos ver que la representación más utilizada es la que tiene $182$ parámetros, y que el resto sigue una distribución algo uniforme.

\begin{figure}[h]
    \centering
    \includegraphics[width=0.8\textwidth]{../graphs/ExperimentsByNInputs.png}
    \caption{Número de experimentos por representación del estado}
    \label{fig:experimentsByNInputs}
\end{figure}

Podemos ver como fue la media de cada entrenamiento para los que tienen el mismo número de parámetros en la ilustración \ref{fig:victoriesByNInputs}.

\begin{figure}[h]
    \centering
    
\includegraphics[width=0.8\textwidth]{images/showdownLogo.png}
    \caption{Porcentaje medio de victorias por representación del estado}
    \label{fig:victoriesByNInputs}
\end{figure}

También podemos ver el porcentaje de victorias entre los experimentos que tienen solo los hiperparámetros de la representación distintos, en la ilustración \ref{fig:victoriesByNInputsAndHyperparameters}.

\begin{figure}[h]
    \centering
    
\includegraphics[width=0.8\textwidth]{images/showdownLogo.png}
    \caption{Matriz de calor en la cual se enseñan para los experimentos en los que solo cambia la representación del estado y los hiperparámetros, el porcentaje de victorias.}
    \label{fig:victoriesByNInputsAndHyperparameters}
\end{figure}

Finalmente también tenemos la ilustración \ref{fig:nInputsBradleyTerry}, que muestra la gráfica del valor Bradley-Terry para los experimentos hechos y la gráfica \ref{fig:nInputsSurrogateBradleyTerry} que muestra los valores de Bradley-Terry mediante el modelo sustitito.

\begin{figure}[h]
    \centering
    
\includegraphics[width=0.8\textwidth]{images/showdownLogo.png}
    \caption{Valores de Bradley-Terry para los experimentos hechos}
    \label{fig:nInputsBradleyTerry}
\end{figure}

\begin{figure}[h]
    \centering
    
\includegraphics[width=0.8\textwidth]{images/showdownLogo.png}
    \caption{Valores de Bradley-Terry para los experimentos del modelo sustituto}
    \label{fig:nInputsSurrogateBradleyTerry}
\end{figure}

Para la manera de representar el campo de batalla, pokémon y movimiento, podemos observar las mismas ilustraciones.

\subsection{Hipótesis \HTwo}

\subsection{Hipótesis \HThree}
